
\subsubsection{6.1 Methodological Contributions}

This work addresses a measurement gap in AI alignment research: existing evaluation measures AI performance without considering reciprocal effects on human oversight behavior. We contribute the first framework to operationalize bidirectional alignment dynamics as coupled, measurable variables (ACE and HOR) with experimental protocols for testing coupling mechanisms.

\textbf{Key advances:}

\begin{enumerate}
\item \textbf{Architectural Separation:} We leverage Constitutional AI's policy-layer architecture to enable testable capability invariance via compliance modulation. The h-e1 gate (ICC > 0.95, ANOVA p > 0.05, Cohen's f < 0.10) enforces this separation as a testable assumption.
\end{enumerate}

\begin{enumerate}
\item \textbf{Integration Across Research Areas:} We unify AI alignment evaluation, automation bias from human factors, and signal detection theory into a single measurement framework.
\end{enumerate}

\begin{enumerate}
\item \textbf{Pre-Registered Falsification Criteria:} Dependency-driven hypothesis decomposition with gate-based validation reduces researcher degrees of freedom. Each hypothesis has pre-specified success thresholds, failure routing logic, and prerequisite dependencies.
\end{enumerate}

\begin{enumerate}
\item \textbf{Automated Quality Control:} The pipeline detected and corrected mock data usage, preventing false validation of prerequisite assumption A1.
\end{enumerate}

\subsubsection{6.2 Limitations}

\textbf{Limitation 1: No Empirical Evidence}

All five predictions (P1-P5) and all five assumptions (A1-A5) remain INCONCLUSIVE/UNVERIFIED. No experiments were executed. Zero empirical data collected on ACE-HOR coupling, automation bias mediation, or metacognitive interventions.

The entire research contribution rests on empirical validation. Without experiment execution, this work contributes only methodological design and engineering, not scientific findings about bidirectional alignment dynamics.

Root cause: API key unavailability blocked h-e1 execution. Expected cost (~$1,620 for full MMLU + HumanEval evaluation) likely required budget approval that was not obtained. Dependent hypotheses (h-m1 through h-m4) could not proceed due to h-e1 prerequisite failure.

Methodological contributions are valuable for emerging research areas. Bidirectional alignment evaluation is a new problem space—measurement frameworks are a necessary first step before large-scale empirical studies. Our work establishes how to measure bidirectional dynamics, enabling future researchers to conduct empirical tests. The infrastructure code is validated and datasets verified, ready for execution when resources become available.

\textbf{Limitation 2: Capability Invariance Assumption Unverified}

Assumption A1 (policy-layer modulation preserves base capability with ICC > 0.95) is the foundational prerequisite for all downstream measurements. If violated, ACE becomes confounded with capability changes, making all ACE-HOR coupling results uninterpretable. This assumption remains empirically untested.

Many alignment interventions do affect capability. Constitutional AI system prompts may alter verbosity, hedging, or reasoning depth in ways that affect MMLU/HumanEval performance. If A1 is violated, the framework's causal inference breaks down.

We explicitly designed h-e1 as a MUST_WORK gate to test this assumption. Failure triggers routing to Phase 0 (architecture re-selection). The assumption is falsifiable and testable, not hidden or implicitly assumed. Constitutional AI architecture suggests capability invariance via policy-layer separation, but this must be validated empirically.

Mitigation: h-e1 execution tests A1 directly. If violated, alternative ACE operationalizations exist: frozen base model with separate trainable policy head, retrieval-augmented constitutional rules, or partial correlation controlling for capability.

\textbf{Limitation 3: Ecological Validity of Seeded Errors}

HOR measurement via signal detection d' requires ground-truth error labels. We use seeded within-distribution errors sampled from the model's empirical confusion matrix (Assumption A2). However, ecological validity—whether seeded errors represent real model failures users encounter in deployment—remains unverified.

If seeded errors are detectability-biased, HOR measurements reflect stress-testing rather than realistic oversight. Results may not generalize to deployment contexts where error types, frequencies, and detectability differ.

Seeded error approach is common in oversight and vigilance research. Sampling from empirical confusion matrix ensures errors resemble model's actual failure modes. We explicitly acknowledge the validity threat.

Mitigation: Compare seeded error HOR measurements with naturally-occurring error detection in deployment logs. If correlation is high (r > 0.7), ecological validity supported.

\textbf{Limitation 4: Generalization Uncertainty}

Even if experiments succeed, the framework's practical utility for real-world deployment remains uncertain. Three concerns: (1) Laboratory seeded errors may not capture the full distribution of deployment failures; (2) Single-session or longitudinal evaluations in controlled settings may not predict oversight quality in real workflow integration with switching costs and multi-tasking; (3) Framework is designed for Claude 3 Opus with Constitutional AI policy-layer manipulation and may not generalize to other architectures with different alignment mechanisms.

Laboratory experiments establish causal mechanisms, to be validated against field data. Mitigation: (1) Correlate lab HOR measurements with deployment error detection rates from production logs; (2) Deploy framework in real workflows and compare lab-predicted vs. field-observed coupling patterns; (3) Test framework on alternative models to assess generalization.

\subsubsection{6.3 Status Summary}

\textbf{Infrastructure Validation:}
\begin{itemize}
\item Quality control effective (mock data detected/fixed, 6/6 verification checks passed)
\item Real datasets verified (MMLU 57 subjects/14,042 questions, HumanEval 164 problems)
\item Dependency structure complete (5 hypotheses with gate criteria, falsification logic)
\item Deviation classification transparent (all IMPLEMENTATION_GAP, no methodology failures)
\end{itemize}

\textbf{Empirical Validation:}
\begin{itemize}
\item No experiments executed (API resource constraint)
\item All predictions (P1-P5) INCONCLUSIVE
\item All assumptions (A1-A5) UNVERIFIED
\end{itemize}

\textbf{Contribution:} Methodological infrastructure (framework design + infrastructure implementation), not empirical findings (experiments pending).
